% !TeX root = ./rpz.tex
\chapter{{\fontsize{14pt}{21pt}\selectfont\MakeUppercase{Анализ предметной области и обоснование выбора технологий}}}
\label{cha:analysis}

\section{Исходное состояние объекта и характеристика предметной области}

В настоящее время большинство пользователей решают задачу планирования
маршрута вручную. Типичный сценарий включает формирование списка дел в
менеджере задач, последующий поиск адресов в навигационном сервисе и
самостоятельное определение порядка посещения точек~\cite{alessandretti2022urban}. При наличии временных
окон (например, «с 9:00 до 12:00») пользователь вынужден дополнительно сопоставлять время перемещения и длительность выполнения задачи.

Существующие цифровые инструменты можно разделить на две группы:
\begin{itemize}
  \item менеджеры задач (не учитывают пространственный фактор);
  \item навигационные сервисы (не работают с временными окнами и списками задач).
\end{itemize}

Таким образом, пользователь вынужден использовать несколько инструментов одновременно,
что увеличивает трудозатраты и повышает вероятность ошибок \cite{reshenie_zadach_planirovaniya2020}.

Целевой аудиторией системы являются:

\begin{itemize}
  \item частные пользователи, планирующие личные дела в городе;
  \item мастера и специалисты на выезде;
  \item курьеры малого бизнеса;
  \item индивидуальные предприниматели;
  \item диспетчеры, формирующие небольшие маршруты вручную.
\end{itemize}

Особенностью данной предметной области является необходимость одновременного учета:
\begin{itemize}
  \item географического расположения объектов;
  \item временных ограничений;
  \item длительности выполнения задач;
  \item минимизации времени перемещения и ожидания.
\end{itemize}

Таким образом, предметная область представляет собой систему с признаками логистической оптимизации, однако в упрощённом масштабе, ориентированном на одного пользователя или малый бизнес.

\section{Постановка проблемы и формирование цели разработки}

Анализ исходного состояния показал, что основной проблемой является отсутствие доступного инструмента, объединяющего управление задачами и автоматическую маршрутизацию с учетом временных окон.

Ручной подход:

\begin{itemize}
  \item требует значительных временных затрат;
  \item не гарантирует оптимальность маршрута;
  \item становится неэффективным при увеличении числа точек;
  \item не обеспечивает формализованного учета временных ограничений.
\end{itemize}

Задача рационального планирования маршрута с учeтом временных ограничений
относится к классу задач маршрутизации с временными окнами
(Vehicle Routing Problem with Time Windows, VRPTW) ---
известной NP-трудной задаче комбинаторной оптимизации \cite{hertrich2019sweep, reshenie_zadach_marshrutizasii2022, iklassov2024svrptw}.

Целью разработки является создание веб-приложения, обеспечивающего
автоматизированное формирование оптимального маршрута выполнения задач
с учетом их географического положения, временных окон и длительности выполнения.

Назначение продукта заключается в цифровой трансформации процесса ручного планирования
маршрута в автоматизированный процесс.

На основании цели сформирован следующий перечень функций:

\begin{itemize}
  \item создание, редактирование и удаление задач;
  \item хранение географических координат;
  \item получение матрицы расстояний через внешние API;
  \item вычисление оптимальной последовательности выполнения задач;
  \item расчет времени прибытия, ожидания и завершения;
  \item визуализация маршрута на карте;
  \item экспорт маршрута в формате JSON.
\end{itemize}

Для обеспечения соответствия назначению определены следующие нефункциональные требования:

\begin{itemize}
  \item время ответа сервера — не более 2 секунд при типовой нагрузке;
  \item поддержка современных браузеров;
  \item масштабируемость до 100 одновременных пользователей;
  \item устойчивость к повторным запросам внешних API за счет кэширования.
\end{itemize}

\section{Обоснование выбора технологий и платформ}

\subsection{Выбор серверной платформы}

В качестве языка для написания серверной части был выбран Golang\cite{golang}.

Выбор языка Golang\cite{golang} обусловлен:

\begin{itemize}
  \item статической типизацией, снижающей риск ошибок;
  \item встроенной моделью конкурентности на основе горутин \cite{saioc2023goroutine};
  \item наличием практического опыта разработки.
\end{itemize}

\subsection{Выбор СУБД и работы с геоданными}

В качестве СУБД выбрана PostgreSQL\cite{PostgreSQL}.

Обоснование выбора:

\begin{itemize}
  \item наличие геометрических индексов (GiST\cite{PostgreSQL_gist});
  \item непрерывное промышленное применение с 1996~года; стабильный релизный цикл и соответствие стандарту SQL:2016~\cite{PostgreSQL};
  \item наличием практического опыта разработки.
\end{itemize}

\subsection{Выбор клиентских технологий}

Для реализации пользовательского интерфейса был выбран React\cite{react}.

Причины выбора:

\begin{itemize}
  \item компонентный подход к разработке;
  \item развитая экосистема;
  \item наличие поддержки картографических библиотек.
\end{itemize}

В проекте применяется TypeScript --- статически типизированное надмножество JavaScript,
что обеспечивает раннее обнаружение ошибок на этапе компиляции.
Эмпирическое исследование 604~проектов на GitHub подтверждает, что TypeScript-проекты
демонстрируют значительно более высокое качество и понятность кода
по сравнению с аналогами на JavaScript~\cite{bogner2022ts}.

\subsection{Выбор картографического API}

Для геокодирования адресов и визуализации маршрута выбран Яндекс Карты JavaScript API~2.1.
В таблице~\ref{tab:maps_compare} представлено сравнение доступных картографических решений.

\begin{table}[H]
  \caption{Сравнение картографических API}
  \label{tab:maps_compare}
  \begin{tabular}{|p{3.5cm}|p{3.5cm}|p{3.5cm}|p{3cm}|}
    \hline
    \centering\textbf{Критерий} & \centering\textbf{Яндекс Карты JS API} & \centering\textbf{Google Maps JS API} & \multicolumn{1}{c|}{\textbf{OpenStreetMap / Leaflet}} \\ \hline
    Геокодер на русском языке   & Да, включая топонимику РФ              & Да                                    & Основан на устарвшей информации                       \\ \hline
    Бесплатный лимит (геокодер) & 1\,000 запр./сутки                     & 200\,\$ кредит/мес.                   & Без ограничений                                       \\ \hline
    Работа в России             & Без ограничений                        & Может быть недоступен                 & Без ограничений                                       \\ \hline
  \end{tabular}
\end{table}

По совокупности критериев таблицы~\ref{tab:maps_compare} для целевой аудитории
(российские пользователи) выбраны Яндекс Карты: единственный из рассмотренных вариантов,
обеспечивающий поддержку российской топонимики, встроенный мультимаршрутизатор для
построения матрицы расстояний и гарантированную доступность в российском сетевом сегменте.

\subsection{Итоговый технологический стек}

В таблице~\ref{tab:tech_stack} приведeн обоснованный итоговый технологический стек.

\begin{table}[H]
  \caption{Итоговый технологический стек разрабатываемой системы}
  \label{tab:tech_stack}
  \begin{tabular}{|p{4cm}|p{3.5cm}|p{6cm}|}
    \hline
    \centering\textbf{Уровень} & \centering\textbf{Технология}           & \multicolumn{1}{c|}{\textbf{Обоснование}}                                           \\ \hline
    Серверная часть            & Go 1.23 + Gin~\cite{gin}                & Нативная компиляция, модель конкурентности goroutine, минимальные накладные расходы \\ \hline
    База данных                & PostgreSQL 16 + PostGIS~\cite{postgis}  & Реляционность, геоиндексы GiST, зрелость платформы                                  \\ \hline
    Клиентская часть           & React 18 + TypeScript~\cite{typescript} & Компонентный подход, статическая типизация, развитая экосистема                     \\ \hline
    Сборка фронтенда           & Vite 6~\cite{vite}                      & Горячая замена модулей, оптимизированная сборка для производственной среды          \\ \hline
    Картографический API       & Яндекс Карты JS API 2.1                 & Русскоязычный геокодер, встроенный мультимаршрутизатор                              \\ \hline
    Контейнеризация            & Docker Compose~\cite{docker}            & Воспроизводимость среды, автоматическое применение миграций                         \\ \hline
  \end{tabular}
\end{table}

\section{Анализ алгоритмических подходов к задаче маршрутизации с временными окнами}

Задача, решаемая в данной работе, относится к классу VRPTW, рассмотренному в предыдущем разделе, ---
обобщению классической задачи коммивояжeра (TSP)~\cite{held1962dp} с дополнительными временными ограничениями.

\subsection{Формальная постановка задачи}

Пусть задан граф $G = (V, E)$, где $V = \{v_0, v_1, \ldots, v_n\}$ --- множество
узлов (задач пользователя), $E$ --- множество дуг с матрицей расстояний $d_{ij}$
и матрицей времeн перемещения $t_{ij}$.

Каждый узел $v_i$ ($i \geq 1$) характеризуется:
\begin{itemize}
  \item временным окном $[a_i,\, b_i]$~--- допустимым интервалом начала обслуживания;
  \item длительностью обслуживания $s_i$ (минуты).
\end{itemize}

Если маршрутный агент прибывает в узел $v_i$ в момент $\tau_i < a_i$, он ожидает
до момента $a_i$; если $\tau_i > b_i - s_i$~--- обслуживание начать невозможно
(нарушение ограничения).

\textbf{Цель:} найти порядок посещения всех узлов, минимизирующий суммарное время
маршрута $T = \sum_{i} t_{ij} + \sum_{i} w_i$, где $w_i = \max(0,\; a_i - \tau_i)$~---
время ожидания в узле $v_i$, при соблюдении условий $\tau_i \leq b_i - s_i$ для всех $i$.

На рисунке~\ref{fig:vrptw_illustration} приведeн пример графа задачи VRPTW
с тремя узлами. Ребра графа подписаны временем перемещения $d_{ij}$ (в минутах),
получаемым из матрицы расстояний внешнего API. Алгоритм NNH-TW выбирает
оптимальный порядок посещения: $v_0 \to v_2 \to v_1$~--- посещение ближайшей
допустимой точки на каждом шаге.

\begin{figure}[H]
  \centering
  \includegraphics[width=0.92\textwidth]{vrptw_illustration.png}
  \caption{Пример графа задачи маршрутизации с временными окнами (VRPTW):
    $d_{ij}$~--- время перемещения между узлами (мин),
    $[a_i,\,b_i]$~--- временное окно, $s_i$~--- длительность обслуживания;
    жирными стрелками выделен оптимальный маршрут $v_0 \to v_2 \to v_1$}
  \label{fig:vrptw_illustration}
\end{figure}

\subsection{Сравнение алгоритмических подходов}

В таблице~\ref{tab:algorithms_compare} приведeн сравнительный анализ основных
подходов к решению VRPTW применительно к задачам данной работы.

\begin{table}[H]
  \caption{Сравнение алгоритмических подходов к задаче VRPTW}
  \label{tab:algorithms_compare}
  \begin{tabular}{|p{3.5cm}|p{2.2cm}|p{2.2cm}|p{5.5cm}|}
    \hline
    \centering\textbf{Подход}                                                         & \centering\textbf{Слож\-ность}               & \centering\textbf{Опти\-маль\-ность} & \multicolumn{1}{c|}{\textbf{Применимость}}             \\ \hline
    Полный перебор (brute force)                                                      & $O(n!)$                                      & Да                                   & Только при $n \leq 10$; неприемлем для веб-приложения  \\ \hline
    Динамическое программирование (алг. Хелда--Карпа~\cite{held1962dp})               & $O(2^n \cdot n^2)$                           & Да                                   & До $n \approx 20$; требует значительного объeма памяти \\ \hline
    Жадная эвристика NNH-TW (ближайший сосед с временными окнами)                     & $O(n^2)$                                     & Нет                                  & Веб-приложение; $n < 50$; предсказуемое время ответа   \\ \hline
    Муравьиный алгоритм (ACO)~\cite{reshenie_zadach_marshrutizasii2022,necula2017aco} & $O(\mathit{iter} \cdot n^2 \cdot m)$         & Нет                                  & Средний масштаб; сложная настройка параметров          \\ \hline
    Генетические алгоритмы (GA)~\cite{gladkov2010ga}                                  & $O(\mathit{gen} \cdot \mathit{pop} \cdot n)$ & Нет                                  & Большой масштаб; длительное время вычислений           \\ \hline
  \end{tabular}
\end{table}

Здесь $n$~--- количество узлов маршрута, $m$~--- количество муравьeв (агентов),
$\mathit{iter}$~--- число итераций, $\mathit{gen}$~--- число поколений,
$\mathit{pop}$~--- размер популяции.

\subsection{Обоснование выбора эвристики NNH-TW}

Для разрабатываемого веб-приложения, ориентированного на частных пользователей
и малый бизнес, был выбран алгоритм <<ближайшего соседа с временными окнами>>
(NNH-TW, Nearest Neighbour Heuristic with Time Windows) по следующим причинам:

\begin{itemize}
  \item \textbf{Квадратичная сложность} $O(n^2)$: при $n = 30$ задачах
        алгоритм выполняет не более $30^2 = 900$ элементарных сравнений,
        что на несколько порядков меньше задержки, вносимой получением
        матрицы расстояний от внешнего API (экспериментально подтверждено
        в~разделе~<<Производительность алгоритма оптимизации>>);
  \item \textbf{Предсказуемость}: время работы линейно зависит от числа задач,
        что гарантирует выполнение нефункционального требования по времени ответа
        (не более 2~секунд);
  \item \textbf{Качество решений}: жадные конструктивные эвристики
        обеспечивают приемлемое качество решений при низких вычислительных
        затратах~\cite{liu2023vrpheuristics}, что достаточно для целей
        планирования личных задач и курьерских маршрутов малого масштаба;
  \item \textbf{Расширяемость}: реализованный интерфейс \texttt{Optimizer}
        позволяет заменить алгоритм на более мощный (например, ACO или GA)
        без изменения остального кода системы.
\end{itemize}

\section{Обзор аналогов и сравнительный анализ}

Существующие программные решения, частично пересекающиеся с задачами
разрабатываемой системы, можно разделить на три класса.

\textbf{Первый класс~--- менеджеры задач}
(Todoist~\cite{todoist}, Microsoft To~Do~\cite{ms_todo}):
реализуют полноценное создание, редактирование и удаление задач,
поддерживают напоминания и многоплатформенную синхронизацию,
однако полностью лишены функций геокодирования адресов и построения маршрутов.

\textbf{Второй класс~--- навигационные сервисы}
(Google Maps~\cite{google_maps}):
обеспечивают прокладку маршрутов по нескольким точкам (не более 10 остановок),
однако не поддерживают автоматическую оптимизацию порядка посещения
и не предоставляют инструментов управления задачами~\cite{google_maps_help}.

\textbf{Третий класс~--- специализированные логистические платформы}
(Routific~\cite{routific}, Relog~\cite{relog},
Яндекс Маршрутизация~\cite{yandex_routing}):
сочетают геокодирование, оптимизацию маршрутов и учeт временных окон,
однако ориентированы на корпоративное управление автопарком
(массовая обработка заказов, несколько водителей, диспетчеры)
и недоступны для индивидуального применения без коммерческой подписки.

В таблице~\ref{tab:analyze_analogs} приведeн сравнительный анализ
представителей всех трeх классов.

\begin{table}[H]
  \caption{Сравнительный анализ аналогов}
  \label{tab:analyze_analogs}
  {\small
    \begin{tabular}{|p{2.6cm}|p{1.6cm}|p{1.7cm}|p{1.7cm}|p{1.7cm}|p{2cm}|p{2.1cm}|}
      \hline
      \centering\textbf{Критерий}
       & \centering\textbf{Todoist}
       & \centering\textbf{Microsoft To Do}
       & \centering\textbf{Google Maps}
       & \centering\textbf{Routific}
       & \centering\textbf{Relog / Яндекс Марш\-рутизация}
       & \multicolumn{1}{p{2.1cm}|}{\centering\textbf{Проекти\-руемое решение}}                              \\ \hline
      Управление задачами
       & Да                                                                     & Да  & Нет & Нет & Нет & Да \\ \hline
      Геокодиро\-вание адресов
       & Нет                                                                    & Нет & Да  & Да  & Да  & Да \\ \hline
      Авто\-мати\-ческая оптимизация маршрута
       & Нет                                                                    & Нет & Нет & Да  & Да  & Да \\ \hline
      Учeт временных окон
       & Нет                                                                    & Нет & Нет & Да  & Нет & Да \\ \hline
      Целевой сегмент
       & Частные пользо\-ватели
       & Частные пользо\-ватели
       & Массовый
       & Корпора\-тивный
       & Средний и крупный бизнес
       & Частный пользователь и малый бизнес                                                                 \\ \hline
      Модель распро\-странения
       & Условно-бесплат\-ная
       & Бесплат\-ный
       & Бесплат\-ный
       & Условно-бесплат\-ная
       & Платная подписка
       & Бесплатный                                                                                          \\ \hline
    \end{tabular}}
\end{table}

Из таблицы~\ref{tab:analyze_analogs} следует, что ни один из рассмотренных аналогов
не обеспечивает одновременно управление задачами, автоматическую оптимизацию
маршрута с учeтом временных окон и доступность для частного пользователя
без коммерческой подписки.
Менеджеры задач первого класса лишены геолокационных функций;
навигационные сервисы второго класса не поддерживают автоматическую
оптимизацию порядка посещения и не управляют задачами;
логистические платформы третьего класса реализуют требуемый функционал,
однако ориентированы на корпоративный автопарк и недоступны
для целевой аудитории разрабатываемой системы.


\section{Выводы}

Проведeн анализ предметной области, выявлены ограничения существующих инструментов
и обоснована актуальность разработки системы планирования маршрутов с учeтом
временных окон. Сформулирована цель и назначение продукта, определeн перечень
функций и нефункциональных требований.

Проведeн сравнительный анализ алгоритмических подходов к задаче VRPTW:
полный перебор, динамическое программирование, жадные эвристики и метаэвристики.
Обоснован выбор алгоритма NNH-TW с квадратичной сложностью $O(n^2)$ как наиболее
подходящего для целевого масштаба системы ($n < 50$) и требований по времени ответа.

На основании сравнительного анализа технологий выбран стек разработки,
обеспечивающий требуемую производительность и масштабируемость системы.
Проведeнный анализ аналогов подтвердил целесообразность создания специализированного
решения, ориентированного на частных пользователей и малый бизнес.
